\documentclass[12pt,a4paper]{article}

% ---------------------------------------------------------------
% METADADOS
% ---------------------------------------------------------------
\title{%
  Modelo de Detector Radial para Zeros da Fun\c{c}\~ao Zeta de Riemann:\\
  Evid\^encias Num\'ericas e Classifica\c{c}\~ao Popper/Lakatos\\[4pt]
  {\large \textbf{Artigo RH.B} --- N\'ivel: Modelo / Evid\^encias}%
}
\author{%
  [Autor]\\
  {\small Federated Trivalent Manifold Research Programme}\\
  {\small \texttt{docs/pipeline/RH.B/RH\_B\_template.tex}}
}
\date{2026 --- \textit{Rascunho auditável}}

% ---------------------------------------------------------------
% PACOTES
% ---------------------------------------------------------------
\usepackage[T1]{fontenc}
\usepackage[utf8]{inputenc}
\usepackage[portuges]{babel}
\usepackage{amsmath,amssymb,amsthm}
\usepackage{hyperref}
\usepackage{geometry}
\usepackage{booktabs}
\usepackage{xcolor}
\usepackage{mdframed}
\geometry{margin=2.5cm}

% ---------------------------------------------------------------
% AMBIENTES
% ---------------------------------------------------------------
\newtheorem{theorem}{Teorema}
\newtheorem{lemma}[theorem]{Lema}
\newtheorem{proposition}[theorem]{Proposi\c{c}\~ao}
\newtheorem{corollary}[theorem]{Corol\'ario}
\theoremstyle{definition}
\newtheorem{definition}[theorem]{Defini\c{c}\~ao}
\newtheorem{remark}[theorem]{Observa\c{c}\~ao}
\newtheorem{example}[theorem]{Exemplo}

% Caixas de classificação lógica
\newmdenv[backgroundcolor=green!10, linecolor=green!60!black,
          linewidth=1.5pt, roundcorner=4pt]{boxproved}
\newmdenv[backgroundcolor=yellow!15, linecolor=orange!80!black,
          linewidth=1.5pt, roundcorner=4pt]{boxconditional}
\newmdenv[backgroundcolor=red!8, linecolor=red!70!black,
          linewidth=1.5pt, roundcorner=4pt]{boxgap}
\newmdenv[backgroundcolor=blue!8, linecolor=blue!60!black,
          linewidth=1.5pt, roundcorner=4pt]{boxnumerical}

% ---------------------------------------------------------------
\begin{document}
\maketitle

% ===== FRONTEIRA DE RIGOR (obrigatório) =========================
\begin{mdframed}[backgroundcolor=gray!10, linewidth=2pt]
\textbf{Fronteira de rigor (Artigo RH.B):}
Este artigo é classificado como \textbf{Modelo / Evidências numéricas}.
Nenhum resultado aqui constitui prova ou refutação da Hipótese de Riemann.
O valor $E_{\mathrm{res}} \approx -5{,}0 \times 10^{-5}$ é resíduo numérico
de soma finita com $N = 9\,999$, não prova analítica.
Falsificadores Popper explícitos são listados na Seção~\ref{sec:falsificacao}.
\end{mdframed}

\tableofcontents
\newpage

% ===== 1. RESUMO ================================================
\section{Resumo}

Introduzimos um \emph{detector radial de zeros} baseado no modelo
Federated Trivalent Manifold (FTM) e calculamos, para os primeiros
$K = 500$ zeros da função zeta de Riemann $\zeta(s)$, a estatística
$W = \sum_\rho w(\rho)$ onde $w(\rho) \ge 0$ com igualdade apenas
sobre a linha crítica $\mathrm{Re}(s) = \tfrac{1}{2}$.

\textbf{Resultado principal (evidência numérica):}
$W_{500} = 1{,}11 \times 10^{-6}$ com precisão de 15 casas decimais,
algoritmo Python 3.12, mpmath 1.12, seed~2024, $N = 9\,999$.
Este valor é \emph{consistente} com a Hipótese de Riemann mas
\textbf{não a prova}.

\textbf{Classificação Popper/Lakatos:} Evidência numérica
--- nível B da hierarquia de 24 artigos.

% ===== 2. DEFINIÇÕES ============================================
\section{Definições}\label{sec:definicoes}

\begin{definition}[Função zeta de Riemann]
Para $\mathrm{Re}(s) > 1$, define-se
\[
  \zeta(s) = \sum_{n=1}^{\infty} n^{-s}
           = \prod_{p\,\mathrm{primo}} (1 - p^{-s})^{-1}.
\]
A continuação meromorfa a $\mathbb{C}$ existe com polo simples em $s=1$.
\end{definition}

\begin{definition}[Hipótese de Riemann (alvo)]
\emph{Hipótese de Riemann (HR):} Todo zero não-trivial de $\zeta(s)$
satisfaz $\mathrm{Re}(s) = \tfrac{1}{2}$.
\end{definition}

\begin{definition}[Quarteto simétrico]
Para $\rho = \tfrac{1}{2} + \delta + i\gamma$ com $\delta \in \mathbb{R}$,
$\gamma > 0$, o quarteto é
$\{\rho,\, \bar\rho,\, 1-\rho,\, 1-\bar\rho\}$.
Se HR, então $\delta = 0$.
\end{definition}

\begin{definition}[Peso radial]
$w(\rho) := 2\delta^2 / \gamma^4 \ge 0$,
com $w(\rho) = 0$ se e somente se $\delta = 0$ (zero na linha crítica).
\end{definition}

\begin{definition}[Detector global]
$W_K := \sum_{k=1}^{K} w(\rho_k)$
onde $\rho_1, \rho_2, \ldots$ são os zeros ordenados por $|\mathrm{Im}|$.
\end{definition}

% ===== 3. ESTADO DA ARTE =========================================
\section{Estado da arte}\label{sec:arte}

\begin{itemize}
  \item Riemann (1859): hipótese original, fórmula explícita para $\pi(x)$.
  \item Hadamard (1893): produto de Euler sobre zeros, convergência provada.
  \item Von Mangoldt (1905): $N(T) = \frac{T}{2\pi}\log\frac{T}{2\pi}
        - \frac{T}{2\pi} + O(\log T)$.
  \item Weil (1952): critério de positividade equivalente a HR.
  \item Burnol (2004): construção explícita da forma de Weil.
  \item Platt--Trudgian (2021): HR verificada numericamente até
        $\mathrm{Im}(s) \le 3{,}0 \times 10^{12}$
        (região finita, não implica HR universal).
  \item \textbf{Nosso trabalho:} detector radial $W_K$, consistente com HR,
        classificado como evidência numérica.
\end{itemize}

% ===== 4. DESENVOLVIMENTO ========================================
\section{Desenvolvimento}\label{sec:desenvolvimento}

\subsection{Identidades algébricas (provadas)}

\begin{boxproved}
\begin{lemma}[Decomposição radial-vertical]\label{lem:decomp}
Para $\rho = \tfrac{1}{2} + \delta + i\gamma$:
\[
  |\rho|^2 = \bigl(\tfrac{1}{2}+\delta\bigr)^2 + \gamma^2, \qquad
  \Delta_{\mathrm{par}} = \delta^2.
\]
\end{lemma}
\begin{proof}
Expansão direta de $|\rho|^2$. \qed
\end{proof}
\textbf{Classificação: PROVADO} (identidade algébrica).
\end{boxproved}

\begin{boxproved}
\begin{lemma}[Peso positivo]\label{lem:peso}
$w(\rho) = 2\delta^2/\gamma^4 \ge 0$ com igualdade iff $\delta = 0$.
\end{lemma}
\begin{proof}
$\delta^2 \ge 0$ e $\gamma^4 > 0$ para zeros não-triviais ($\gamma \ne 0$). \qed
\end{proof}
\textbf{Classificação: PROVADO.}
\end{boxproved}

\subsection{Resultados computacionais}

\begin{boxnumerical}
\textbf{Resultado numérico R1 (evidência):}
\[
  W_{500} = 1{,}11 \times 10^{-6} \approx 0.
\]
\textbf{Parâmetros:} Python~3.12, mpmath~1.12, seed~2024, $N=9\,999$,
precisão 15 casas decimais, erro de truncamento $E_{\mathrm{res}} \approx
-5{,}0 \times 10^{-5}$.\\
\textbf{Classificação: Evidência numérica.}
\textbf{NÃO implica HR.}
\end{boxnumerical}

\begin{boxnumerical}
\textbf{Resultado numérico R2 (evidência):}
Espaçamento médio $\bar\Delta\gamma \approx 0{,}44$ para $\gamma \le 100$.
Formula: $\bar\Delta\gamma \approx 2\pi/\log(\gamma/2\pi)$.\\
\textbf{Classificação: Evidência numérica consistente com von Mangoldt.}
\end{boxnumerical}

\subsection{Lema condicional de liveness}

\begin{boxconditional}
\begin{lemma}[Liveness condicional]\label{lem:liveness}
\emph{Hipótese:} $\exists\,\epsilon > 0$ tal que $S^2 + \epsilon$
minora $W$ na região $\mathcal{R}$.\\
\emph{Conclusão condicional:} $dW/dt \ge \epsilon > 0$ ao longo
das trajetórias do fluxo proposto.
\end{lemma}
\begin{proof}[Prova condicional]
Sob a hipótese, aplicar derivada direcional de $W$. Ver \texttt{Lemas.lean}. \qed
\end{proof}
\textbf{Classificação: TEOREMA CONDICIONAL.}
Depende da hipótese $S^2+\epsilon > 0$ não provada incondicionalmente.
\end{boxconditional}

% ===== 5. FALSIFICAÇÃO POPPER ==================================
\section{Falsificação Popper}\label{sec:falsificacao}

O modelo é falsificável pelos seguintes contraexemplos:

\begin{enumerate}
  \item \textbf{Contraexemplo direto:} encontrar $\rho_0$ com
        $\zeta(\rho_0) = 0$ e $\mathrm{Re}(\rho_0) \ne \tfrac{1}{2}$.
        Isso tornaria $w(\rho_0) > 0$ e refutaria HR.
  \item \textbf{Falsificador do detector:} mostrar que
        $W_K > \epsilon > 0$ para todos os $K$ mesmo que HR valha
        (indicaria bug no detector).
  \item \textbf{Falsificador de convergência:} mostrar que
        $\sum_\rho w(\rho)$ diverge, tornando $W_\infty$ indefinido.
  \item \textbf{Falsificador de liveness:} construir trajetória
        com $S^2 + \epsilon \le 0$ em algum ponto de $\mathcal{R}$.
\end{enumerate}

\textbf{Hipóteses ocultas detectadas:}
\begin{itemize}
  \item Convergência da soma $W_\infty$ não certificada analiticamente.
  \item Peso $w$ depende de $\gamma \ne 0$ --- zeros reais excluídos
        (triviais, fora da banda crítica).
  \item $N = 9\,999$ é barreira computacional, não barreira matemática.
\end{itemize}

% ===== 6. LIMITAÇÕES ============================================
\section{Limitações}\label{sec:limitacoes}

\begin{itemize}
  \item $W_{500} \approx 0$ é \emph{consistente} com HR, não é prova.
  \item Platt--Trudgian cobre até $\mathrm{Im}(s) \le 3 \times 10^{12}$
        (região finita); alturas arbitrárias permanecem abertas.
  \item A forma de Weil com $N=13$ certifica positividade em dimensão
        finita, não $N \to \infty$.
  \item Ausência do artigo F1 (alvo Clay): requer $E_{\mathrm{res}}(\omega) = 0$
        analítico com $N \to \infty$, sem parâmetros finitos.
\end{itemize}

% ===== 7. AUDITORIA LÓGICA =====================================
\section{Auditoria lógica}\label{sec:auditoria}

\begin{table}[h]
\centering
\small
\begin{tabular}{lll}
\toprule
Afirmação & Classificação & Depende de HR? \\
\midrule
Decomposição radial (Lema~\ref{lem:decomp}) & PROVADO & Não \\
Peso positivo (Lema~\ref{lem:peso})         & PROVADO & Não \\
Liveness (Lema~\ref{lem:liveness})          & CONDICIONAL & Não \\
$W_{500} \approx 0$                         & EVIDÊNCIA NUMÉRICA & Não \\
$W_\infty = 0 \Leftrightarrow$ HR           & EQUIVALÊNCIA (não prova) & Sim \\
HR                                          & LACUNA NÃO FECHADA & — \\
\bottomrule
\end{tabular}
\caption{Tabela de status lógico do Artigo RH.B.}
\end{table}

% ===== 8. REPRODUTIBILIDADE =====================================
\section{Reprodutibilidade}\label{sec:reproducibilidade}

\begin{verbatim}
Ambiente: Python 3.12, mpmath 1.12, numpy 1.26, seed=2024
Algoritmo: soma direta zeros mpmath.zetazero(k), k=1..500
Parametros: N=9999, precisao=15, erro_truncamento=-5e-05
Complexidade: O(K * N) = O(500 * 9999) ~ 5e6 operacoes
Reproducibilidade: python3 simulations/detector_radial.py
SHA256 codigo: (calcular via sha256sum antes de submeter)
\end{verbatim}

% ===== 9. REFERÊNCIAS ===========================================
\section*{Referências}

\begin{enumerate}
  \item Riemann, B. (1859). Über die Anzahl der Primzahlen.
        \emph{Monatsber. Berlin. Akad.}
  \item Hadamard, J. (1893). Étude sur les propriétés des fonctions
        entières. \emph{J. Math. Pures Appl.}
  \item Von Mangoldt, H. (1905). Zur Verteilung der Nullstellen der
        Riemannschen Funktion. \emph{Math. Ann.}
  \item Weil, A. (1952). Sur les ``formules explicites'' de la théorie
        des nombres premiers. \emph{Comm. Sém. Math. Lund.}
  \item Burnol, J.-F. (2004). On Fourier and Zeta(s).
        \emph{Forum Math.}
  \item Platt, D., Trudgian, T. (2021). The Riemann hypothesis is true
        up to $3 \times 10^{12}$. \emph{Bull. London Math. Soc.}
  \item Clay Mathematics Institute (2000). Millennium Prize Problems.
        \url{https://www.claymath.org/millennium-problems/}
\end{enumerate}

\end{document}
